\section{Conclusion}\label{sec:conclusion}
The useful unit of defense evidence is an evaluated result tied to a receiving view, an execution boundary and a workload. Organizing the literature at this level yields three actionable decisions: protect the final emitter, validate cover where it executes, and use adaptation effects and costs for the mechanisms actually measured. WRA's privacy benefit, AHB's state effects and the cache components' engineering measurements each retain their value without being forced into one ranking.

The contribution is the bridge from existing measurements to deployment choices. The source-linked cases identify when a result can be reused; the comparability audit records which proposed inferences are supported, mismatched or still need a binding. Controlled intervention confirms the practical distinction: upstream size padding can end before a new emitter, and task-preserving cover depends on its execution route. These positive and negative cases turn a diverse side-channel literature into evidence a system designer can apply and a reviewer can reconstruct.
